\chapter{Conclusion}
\label{ch:conclusion}

\section{Summary}
This thesis trains deep-learning and classical time series classifiers to detect and classify bifurcation-driven critical transitions, following the approach \citet{bury2021} introduced: train on labeled synthetic data where the ground truth is known by construction, then test on real data where it is not (Chapter~\ref{ch:data}). The real test case is Mediterranean sapropel formation -- an anoxia transition recorded in three sediment cores -- using the redox-specific Mo and U proxies rather than the full raw element set (Chapter~\ref{ch:data}, Section~\ref{sec:pangaea_provenance}). Unlike \citet{bury2021}'s single CNN-LSTM classifier, this thesis benchmarks 7 deep-learning architectures and 22 classical time series classifiers on the same protocol, at two sequence lengths (Chapter~\ref{ch:methodology}). Chapter~\ref{ch:results} reports what that comparison found, on both the labeled synthetic data, where accuracy is checkable, and the unlabeled real cores, where it is not.

\section{What This Thesis Found}

\subsection{Detection: A Real, Robust Signal}
Cross-referencing the labeled synthetic results against the empirical real-data results separates the 29-model roster into three groups, not a single ranked list (Section~\ref{sec:three_way}): 15 models learned the labeled task and the same skill transfers to the real cores; one (\texttt{boss}) learned the task but does not transfer; the remaining nine never reliably learned it on either dataset. Averaging the 15 transferring models together gives a combined detector reaching a mean AUC of 0.964 across the 14 real segments this thesis evaluates (Section~\ref{sec:ensemble_result}) -- evidence that these classifiers are picking up a genuine pre-transition signal in the real geochemical record, not fitting to noise, since an ensemble of independently-trained architectures cannot reach that score by averaging noise together. This is the central positive result of this thesis: detection, specifically, is real and robust across a much larger model comparison than has been run on this task before.

\subsection{Bifurcation Type: Not Reliably Determined}
Whether a given sapropel was produced by a fold, transcritical, or Hopf bifurcation is a separate question from whether a transition was detected, and this thesis's data cannot answer it reliably. Chance-corrected agreement between the 15 trustworthy models' type votes is indistinguishable from random guessing (Fleiss' $\kappa = 0.007$), and only one of the seven real sapropels clears a bootstrap-confidence threshold for a definitive verdict (Section~\ref{sec:type_reliability}). Where the models do lean toward a type, it is toward Hopf, not the fold type a sapropel onset is normally expected to be: this thesis's pooled favoured-type split is 21\% fold / 59\% hopf / 15\% transcritical / 5\% null, close to a mirror image of \citet{bury2021}'s own published 67\% fold / 17\% hopf / 16\% transcritical / 0\% null on the same anoxia data. \citet{babazadeh2025} report the same qualitative mismatch from a single classifier on an overlapping subset of the same real cores; this thesis's version is built from fifteen independently-trained architectures rather than one, which is a stronger basis for treating the mismatch as a real property of the problem rather than one model's idiosyncrasy -- but Table~\ref{tab:type_reliability} shows even these fifteen models are not reliably agreeing with each other on which non-fold type it actually is, so this remains an open question this thesis surfaces rather than resolves.

\section{How This Compares to Prior Work}
Table~\ref{tab:related_work_comparison} compares this thesis directly against \citet{bury2021} and \citet{babazadeh2025}, the two other deep-learning approaches to this specific problem. The comparison this thesis adds that neither of the other two report is the model-count scale (29 architectures against one each) and the direct, quantified measurement of cross-model type-diagnosis reliability (Section~\ref{sec:type_reliability}) -- without a broad comparison, \citet{bury2021}'s single classifier's type vote has no way to be checked against independent architectures, and this thesis's finding that even fifteen well-performing, independently-trained models do not agree with each other is only visible because that comparison was run. The sapropel selection this thesis tests on is an exact match to \citet{ma2025sdml}'s own seven-sapropel set (Section~\ref{sec:pangaea_provenance}), not \citet{bury2021}'s larger thirteen, which keeps this thesis's real-data numbers comparable to the most directly relevant existing empirical baseline on this exact data.

\section{Limitations}
Three limitations are worth stating plainly, each already established in an earlier chapter and gathered here rather than restated in full.

\begin{itemize}
    \item \textbf{The 29-model roster is not entirely complete.} 17 models have a fully complete synthetic and empirical evaluation at every length they were trained at; the other 12 have at least one gap, ranging from a training run that never finished (\texttt{tde}, \texttt{pf}) to a single missing evaluation length (Section~\ref{sec:roster_status}). Every headline number in Chapter~\ref{ch:results} only ever uses a complete evaluation for a given model and length, so this does not affect the results reported, but it does mean the roster comparison is not yet the full 29-model picture it will eventually be.
    \item \textbf{Noise robustness was not tested.} Every model in this thesis was trained and evaluated on white-noise synthetic data only (Table~\ref{tab:related_work_comparison}); \citet{babazadeh2025}, by contrast, explicitly test coloured noise as well as white. Whether this thesis's roster holds up under a different noise structure is untested.
    \item \textbf{The real test case is one system.} All empirical results in this thesis come from Mediterranean sapropel formation specifically. \citet{bury2021} additionally test paleoclimate transitions and a thermoacoustic rig; this thesis's single real system is a deliberate narrowing to get a larger, corrected version of one case (Section~\ref{sec:label_correction}) rather than a shallower pass over several, but it means the detection result in Section~\ref{sec:ensemble_result} is not yet shown to generalise beyond anoxia.
\end{itemize}

\section{Future Work}
Three extensions follow directly from the limitations above and from gaps this thesis found but did not close. Finishing the roster (retraining \texttt{tde}, \texttt{pf}, and re-running the partial evaluations named in Section~\ref{sec:roster_status}) would complete the 29-model comparison. Testing this thesis's roster against coloured-noise synthetic data, following \citet{babazadeh2025}'s protocol, would show whether the detection result in Section~\ref{sec:ensemble_result} is specific to white-noise training. And the Hopf-versus-fold disagreement in Section~\ref{sec:type_reliability} is, on its own, a specific open question worth its own follow-up: whether it reflects something genuine about this real anoxia record that \citet{bury2021}'s single classifier and this thesis's fifteen happen to read differently, or a property of the type-classification task itself that a broader model comparison on \emph{other} real transitions would help settle.
